\subsection{张量积的复习}

设 A 是环，E 是右 A-模，M 是左 A-模。在《代数学》第二章 §3 第 1 目中我们定义了张量积 \(\mathbf{E} \otimes_{\mathbf{A}} \mathbf{M}\) ，它是一个 \(\mathbf{Z}\) -模。若 \(\mathbf{E}'\) （相应地 \(\mathbf{M}'\) ）是右（相应地左）A-模，而 \(u: \mathbf{E} \to \mathbf{E}'\) （相应地 \(v: \mathbf{M} \to \mathbf{M}'\) ）是同态，我们也（出处同上，第 2 目）定义了一个 \(\mathbf{Z}\) -同态

\[
\boldsymbol {u} \otimes v: \mathrm{E} \otimes_ {\mathrm{A}} \mathrm{M} \rightarrow \mathrm{E} ^ {\prime} \otimes_ {\mathrm{A}} \mathrm{M} ^ {\prime}.
\]

引理 1. 设 \(\mathbf{M}' \xrightarrow{v} \mathbf{M} \xrightarrow{w} \mathbf{M}'' \to 0\) 是左 A-模的一个正合列，E 是一个右 A-模。则序列

\[
\mathrm{E} \otimes_ {\mathrm{A}} \mathrm{M} ^ {\prime} \xrightarrow {1 \otimes v} \mathrm{E} \otimes_ {\mathrm{A}} \mathrm{M} \xrightarrow {1 \otimes w} \mathrm{E} \otimes_ {\mathrm{A}} \mathrm{M} ^ {\prime \prime} \longrightarrow 0
\]

是交换群的正合列。

这就是《代数学》第二章 §3 第 6 目命题 5 的推论。

由此可知，对任意左 A-模同态 \(u: M \rightarrow N\) ，

\[
\mathrm{E} \otimes_ {\mathbf {A}} (\text { Coker } u)
\]

典范地等同于 \(\operatorname{Coker}(1_{E} \otimes u)\) ，这由把引理 1 应用于正合列

\[
\mathrm{M} \xrightarrow {u} \mathrm{N} \longrightarrow \text { Coker } u \longrightarrow 0.
\]

即得。在引理 1 的记号下，我们知道（出处同上），若 v 是单射，即若序列 \(0 \rightarrow M' \stackrel{v}{\rightarrow} M \stackrel{w}{\rightarrow} M'' \rightarrow 0\) 正合，并不能由此推出 \(1_{E} \otimes v\) 是单射，因此一般不能把 \(E \otimes_{A} M'\) 等同于 E \(\otimes_{\mathbf{A}}\) M 的一个子群。然而回顾（《代数学》第二章 §3 第 7 目命题 7 的推论 5）下面的结果：

引理 2. 若 \(v: \mathbf{M}' \to \mathbf{M}\) 是单射且 \(v(\mathbf{M}')\) 是 \(\mathbf{M}\) 的一个直因子，则同态 \(1_{\mathbf{E}} \otimes v\) 是单射，且它的像是 \(\mathbf{E} \otimes_{\mathbf{A}} \mathbf{M}\) 的一个直因子。

